\uuid{u3KS}
\titre{Notion d'équation différentielle}
\auteur{bodin ; croizat ; sacre}
\organisation{exo7}
\licence{CC BY-NC-SA 4.0}
\datecreate{2022-03-08}
\niveau{L1}
\module{Analyse}
\chapitre{Equation différentielle}
\sousChapitre{Résolution d'équation différentielle du premier ordre}
\difficulte{2}
\theme{Notion d'équation différentielle}
\exo7qcmid{hilisit.equations-differentielles.020}

\contenu{
\begin{question}
\qclassification{225.01}
Soit l'équation différentielle $y'- 2xy = 4x$.
Quelles sont les affirmations vraies concernant les solutions de cette équation ?
\begin{answers}
  \good{$y = -2$ est une solution.}
  \bad{$y = +2$ est une solution.}
  \bad{$y = e^{x^2}+2$ est une solution.}
  \good{$y = e^{x^2}-2$ est une solution.}
\end{answers}
\begin{explanations}
Si $y=C$ est constante, alors $y'=0$ et on a $0-2x \cdot C = 4x$ donc $C=-2$ est la seule solution constante de notre équation différentielle. D'autre part, la dérivée de \(e^{x^2}\) étant \(2x e^{x^2}\), on vérifie en remplaçant dans l'équation différentielle que \(e^{x^2}-2\) est solution puisqu'alors $y'-2xy = 2xe^{x^2}-2x(e^{x^2}-2)=4x$. En revanche \(e^{x^2}+2\) n'est pas solution puisque $y'-2xy = 2xe^{x^2}-2x(e^{x^2}+2) = -4x$.
\end{explanations}
\end{question}
}
